\chapter{Related Work}
\label{chap:related}

Automated protein function prediction has been an active research area for more
than two decades. Methods can be grouped into four broad families: sequence
alignment transfer, profile and domain matching, supervised machine learning, and
embedding-based approaches. This chapter surveys each family, identifies their
limitations, and contextualises PROTEA within the landscape.

% ────────────────────────────────────────────────────────────
\section{Alignment-Based Annotation Transfer}
\label{sec:related-alignment}

The foundational insight of alignment-based methods is that sequence similarity
implies functional similarity. If a query protein shares high sequence identity
with a characterised protein, one can transfer the known annotations to the query.

\textbf{BLAST}~\cite{altschul1990blast} remains the most widely deployed tool for
this purpose. It uses a heuristic seed-and-extend strategy to find local
alignments in sublinear time. \texttt{BLASTP} (protein--protein) identifies
homologues in Swiss-Prot and transfers their GO terms to the query. This simple
approach is accurate when identity exceeds roughly 50\,\% and degrades
progressively as identity falls. Below the ``twilight zone'' of $\approx$30\,\%
identity, alignment scores lose statistical power and functional inference becomes
unreliable~\cite{rost1999twilight}.

\textbf{DIAMOND}~\cite{buchfels2021diamond} offers BLAST-comparable sensitivity
at orders-of-magnitude higher throughput, making it practical for genome-scale
annotation. Its speed advantage comes from reduced-alphabet seeds and better
SIMD utilisation; the fundamental twilight-zone limitation remains.

\textbf{Annotation propagation rules} add another layer of nuance: not all GO
terms are equally transferable. Terms annotated with \texttt{IEA} codes can be
propagated freely, while experimental annotations require explicit curation
policies. The CAFA community has established benchmark protocols that penalise
propagation errors~\cite{zhou2019cafa3}.

% ────────────────────────────────────────────────────────────
\section{Profile and Domain Methods}
\label{sec:related-hmm}

To extend sensitivity below the twilight zone, profile methods compare a query
against position-specific models built from multiple sequence alignments of
protein families.

\textbf{HMMER}~\cite{eddy2011hmmer} implements profile \glspl{hmm} for this
purpose. A family is represented as a probabilistic model over residue
frequencies at each position, capturing conserved and variable sites more
faithfully than a single representative sequence. HMMER achieves greater
sensitivity than BLAST at the cost of requiring pre-built profiles.

\textbf{InterPro}~\cite{blum2021interpro} integrates over a dozen family and
domain databases — Pfam, PRINTS, ProSite, TIGRFAM, and others — into a unified
annotation resource. A protein is assigned InterPro entries whenever it matches
any member database; GO terms are then inferred via curated mappings. InterPro
annotation is one of the most reliable computational methods available, but it
is limited to protein families and domains that have been explicitly curated and
modelled.

\textbf{Cascade homology search}~\cite{park1998sequence} iteratively extends
sensitivity by feeding BLAST hits into further BLAST searches (PSI-BLAST) or
HMM construction, at the cost of introducing false positives in remote homology
regimes.

% ────────────────────────────────────────────────────────────
\section{Supervised Machine Learning Approaches}
\label{sec:related-ml}

Supervised methods learn classifiers — one per GO term of interest — from
labelled protein--function pairs. Early approaches used \glspl{hmm} or support
vector machines on curated feature sets (amino acid composition, physicochemical
properties, secondary structure predictions). More recent work uses deep neural
networks applied directly to sequences.

\textbf{DeepGO}~\cite{kulmanov2018deepgo} and its successor
\textbf{DeepGOPlus}~\cite{kulmanov2020deepgoplus} train convolutional networks on
sequence $k$-mers and combine the output with sequence similarity scores. They
demonstrate that learned sequence features and alignment evidence are
complementary: the two signals fuse into a more reliable predictor than either
alone.

\textbf{NetGO}~\cite{you2019netgo} augments sequence-based models with
protein–protein interaction network features, achieving state-of-the-art results
in the CAFA3 challenge. Its dependency on interaction data limits applicability
for poorly studied organisms.

A fundamental challenge for supervised methods is the extreme imbalance and
hierarchy of the GO label space. The number of proteins annotated with any
specific GO term decreases rapidly as one descends the ontology DAG; for many
leaf terms, fewer than ten training examples exist. Hierarchical loss functions
and ontology-aware evaluation metrics (Fmax, AUPR) have been developed to address
this~\cite{clark2013information}.

% ────────────────────────────────────────────────────────────
\section{Protein Language Model Embeddings}
\label{sec:related-plm}

The advent of large-scale self-supervised \glspl{plm} has opened a qualitatively
different avenue for function prediction. Models trained on hundreds of millions
of sequences develop internal representations — embeddings — that capture
evolutionary relationships, secondary and tertiary structure propensities, and
functional sites, without any explicit structural or functional supervision.

\textbf{ESM-1b} and \textbf{ESM-2}~\cite{lin2023esm2} (Meta AI) are transformer
models trained on UniRef50/90. ESM-2 scales from 8M to 15B parameters; the
larger variants generate embeddings that outperform BLAST-based transfer at
low-identity regimes. ESM-2 embeddings serve as the primary representation in
PROTEA.

\textbf{ProtTrans}~\cite{elnaggar2022prottrans} provides a suite of encoder
models (ProtBERT, ProtXLNet, ProtT5) trained on BFD and UniRef100 using
distributed HPC resources. ProtT5-XL achieves the best per-residue and per-protein
accuracy among the ProtTrans variants.

\textbf{Annotation transfer via nearest neighbours} is the inference strategy
most naturally paired with embeddings. Given a set of reference proteins with
known annotations and their embeddings, the $k$ nearest reference neighbours of
a query (by cosine or Euclidean distance) are retrieved and their annotations
aggregated to form a prediction. This approach was explored in
\textcite{littmann2021embeddings} and forms the core inference mechanism of
PROTEA.

\textbf{Scalability} is the central engineering challenge for embedding-based
transfer. Swiss-Prot alone contains over 570,000 entries; computing exact cosine
distances against all of them for a large query set requires $O(NQ)$ dot products,
which becomes infeasible at scale. Approximate nearest-neighbour libraries such
as FAISS~\cite{johnson2021faiss} reduce this to $O(Q \sqrt{N})$ with minimal
accuracy loss for IVFFlat indices.

% ────────────────────────────────────────────────────────────
\section{CAFA Benchmark and Evaluation Standards}
\label{sec:related-cafa}

The Critical Assessment of Functional Annotation (CAFA)~\cite{zhou2019cafa3}
challenge provides the community standard for evaluating function prediction
methods. Participants submit predictions for a set of target proteins; ground
truth is established retrospectively from experimental annotations deposited
after the submission deadline.

The primary metric is $F_{\max}$: the maximum $F_1$ score across all prediction
confidence thresholds, computed separately for each GO aspect. AUPR (area under
the precision--recall curve) and Smin (semantic distance) are complementary
measures. PROTEA's evaluation in \cref{chap:evaluation} follows CAFA conventions.

% ────────────────────────────────────────────────────────────
\section{Annotation Platforms and Pipelines}
\label{sec:related-platforms}

Several platforms have been developed to operationalise function prediction at
genome scale.

\textbf{Argot2}~\cite{falda2012argot2} combines BLAST and HMMER hits with a
graph propagation scheme over the GO DAG. It produces confidence scores by
integrating multiple lines of evidence.

\textbf{PANNZER2}~\cite{toronen2018pannzer2} is a web server that combines
BLAST-based nearest-neighbour transfer with a SANS scoring model. It provides
an interactive interface and supports batch submission.

\textbf{eggNOG-mapper}~\cite{cantalapiedra2021eggnogmapper} maps sequences
against clusters of orthologous groups (COGs) and transfers functional
annotations from the eggNOG database. It is widely used in comparative genomics
pipelines.

None of these platforms provides a reproducible, versioned annotation pipeline
backed by a relational data model that records which ontology version, reference
annotation set, and embedding configuration produced each prediction — a gap
that PROTEA addresses directly.

% ────────────────────────────────────────────────────────────
\section{Summary and Positioning of PROTEA}
\label{sec:related-summary}

\Cref{tab:related-comparison} summarises the methods discussed above along four
axes relevant to this thesis.

\begin{table}[htbp]
  \centering
  \caption{Comparison of protein function prediction approaches.}
  \label{tab:related-comparison}
  \resizebox{\textwidth}{!}{%
  \begin{tabular}{lcccc}
    \toprule
    \textbf{Method} & \textbf{Low-identity} & \textbf{Scalable} & \textbf{Versioned} & \textbf{Open source} \\
    \midrule
    BLAST transfer      & \texttimes & \checkmark & \texttimes & \checkmark \\
    HMMER/InterPro      & \checkmark & \checkmark & \texttimes & \checkmark \\
    DeepGO/DeepGOPlus   & \checkmark & \checkmark & \texttimes & \checkmark \\
    PANNZER2            & \texttimes & \checkmark & \texttimes & \checkmark \\
    ESM-2 kNN transfer  & \checkmark & $\sim$     & \texttimes & \checkmark \\
    \textbf{PROTEA}     & \checkmark & \checkmark & \checkmark & \checkmark \\
    \bottomrule
  \end{tabular}
  }
\end{table}

PROTEA occupies a niche that none of the surveyed systems fills: an end-to-end
distributed pipeline that combines \gls{plm} embeddings with optional alignment
and taxonomy features, records all versioning information in a relational schema,
and exposes the full pipeline through a REST API and web interface.
